% Transcription — fonds Alexandre Grothendieck, shelfmark 31, batch 4 (pages 61-80).
% Established from the facsimile archives/batches/31/batch-04.pdf (a mirror of
% https://grothendieck.umontpellier.fr/31.pdf), per the skill
% transcribe-grothendieck.
%
% Pass: Opus 5 (claude-opus-5), 2026-09-11 — first pass, unchecked against
% the pages by a human.
%
% The batch holds the end of one run, a long run opened by a cover sheet in his
% hand, a corrected typescript leaf, and a second long run.
%
%   Page 61 — ink: the last implication of the Proposition begun on page 59
%     (batch 3), (ii) ⟹ (i bis), through a flasque resolution.
%   Page 62 — an otherwise blank sheet inscribed « cd des affines » in his
%     hand; it opens the run of pages 64-70 and takes no \page{}.
%   Pages 64-70 — ink, mostly fair: ℓ-elementary sheaves i_*(G)_{U,X} on a
%     noetherian X, the Proposition that every noetherian ℓ-torsion sheaf has
%     a finite composition series with strictly elementary quotients, then
%     Cor 1 / Cor 1 bis (criteria for cd_ℓ(X) ≤ n and for the vanishing of
%     H^i_T), Cor 2 / Cor 2 bis (the same with a function φ in place of the
%     integer n), the Proposition of page 66 with its conditions a) and b) on
%     φ, the proof on pages 66-67 by noetherian induction through g_*(G) and
%     the spectral sequence H^p(X, R^q g_*(G)), the note of page 68 bounding
%     φ(x) by φ(y) + 2 deg tr k(x)/k(y), the Corollaire of page 69 with the
%     boxed φ(x) = sup_{y∈x̄}(cd_ℓ k(y) + 2 dim O_{x̄,y}) and its four
%     examples, the marginal Question on cd_ℓ K ≥ dim A + cd_ℓ k, and the
%     Lemme boxed on page 70.
%   Page 71 — a leaf of typescript, scanned upside down, corrected in his hand
%     (references filled in, « VII » made « VIII », a parenthesis struck):
%     2.3 on quotients by groups of multiplicative type, Corollaire 2.4 on the
%     graph morphism, Corollaire 2.5 on a monomorphism of S-group schemes.
%     Insertions in \add{}, deletions in \struck{}.
%   Pages 73-80 — ink, a fast draft: Théorème 1_n on a strictly local
%     noetherian ring of dimension n and a sheaf null in codimension < m,
%     H^i(X,F) = 0 for i > n+1-m; Théorème 2_n, cd_ℓ(X_f) ≤ n; three lemmas
%     relating 1'_n and 2'_n, the reduction to affine space E^d_k, and the
%     final Théorème bounding cd_ℓ(X) by Sup[2 dim X_0, 2 dim X' + 1] for X of
%     finite type over a strictly local A, closing on the embedding
%     X ⊂ E = Y[t_1,…,t_r] and the induction on r.
%
% Pages 63 and 72 (blank) are absent.
%
% The connective prose of pages 73-80 is written at speed and is largely lost;
% the statements and the formulas carry. The apparatus records that honestly
% rather than filling it in.
\documentclass[11pt,a4paper]{article}
% Le préambule commun à toutes les transcriptions du fonds Grothendieck.
%
% Il tient en une idée : **l'incertitude se compose, elle ne se résout pas.**
% Une transcription qui lisse un mot douteux a détruit la seule chose pour
% laquelle on la faisait — un lecteur doit pouvoir savoir, à chaque endroit,
% ce qui était sur le feuillet et ce qui a été ajouté. Les six macros qui
% suivent sont donc l'appareil critique, et non de la décoration.
%
% Les mêmes commandes sont reconnues par `scripts/render.mjs`, qui produit la
% vue de lecture du site. Une macro ajoutée ici doit l'être aussi là-bas,
% faute de quoi le rendu échouera bruyamment — ce qui est voulu : un rendu
% silencieusement faux est pire qu'un rendu absent.

\ProvidesPackage{grothendieck}[2026/08/08 Transcriptions du fonds Grothendieck]

\usepackage{amsmath,amssymb,amsthm}
\usepackage{mathrsfs}
\usepackage{stmaryrd}
% Les diagrammes commutatifs sont partout dans ce fonds : c'est la langue
% même de la géométrie algébrique de Grothendieck, et une transcription qui
% les rendrait en prose ne transcrirait rien.
\usepackage{tikz-cd}
% Français par défaut — c'est la langue des feuillets — mais l'anglais est
% chargé aussi : la traduction et le résumé sont des documents anglais, et la
% typographie française (espaces avant les deux-points) y serait une faute.
% Ces documents basculent par \selectlanguage{english} en tête de corps.
\usepackage[english,french]{babel}
\usepackage{fontspec}
\usepackage{geometry}
\geometry{a4paper,margin=25mm,marginparwidth=28mm}
\usepackage{marginnote}
\usepackage{xcolor}
% ulem, not soul. `soul`'s \st and \ul are built on a character-by-character
% scan that XeLaTeX's font machinery breaks: the document fails with an
% undefined control sequence rather than degrading. ulem does the same two jobs
% and is Unicode-engine safe. `normalem` stops it hijacking \emph, which the
% transcriptions use for Grothendieck's own emphasis.
\usepackage[normalem]{ulem}
\usepackage{hyperref}
\hypersetup{colorlinks=true,linkcolor=black,urlcolor=black,pdfborder={0 0 0}}

% --- Métadonnées du lot -----------------------------------------------------
% Reprises telles quelles par le script de rendu, qui les affiche en tête de
% la vue de lecture. Elles fixent ce que le fichier prétend couvrir : sans
% elles, une transcription flotte sans point d'attache dans un fonds de seize
% mille feuillets.

\newcommand{\folder}[1]{\gdef\grothFolder{#1}}
\newcommand{\batch}[1]{\gdef\grothBatch{#1}}
\newcommand{\leaves}[2]{\gdef\grothFirst{#1}\gdef\grothLast{#2}}
\newcommand{\foldertitle}[1]{\gdef\grothFoldertitle{#1}}
\gdef\grothFolder{?}\gdef\grothBatch{?}\gdef\grothFirst{?}\gdef\grothLast{?}\gdef\grothFoldertitle{}

% --- Le résumé --------------------------------------------------------------
%
% Il ouvre la lecture modernisée, sous le titre « Résumé », et il est composé
% en petit. Ce n'est pas un ornement : c'est le seul passage du document qui
% ne soit pas des mathématiques, et un lecteur doit pouvoir le reconnaître
% avant de l'avoir lu — puis le sauter s'il connaît déjà le sujet, ou s'y
% arrêter s'il ne le connaît pas. Le filet qui le suit marque la reprise du
% corps.

\newenvironment{resume}
  {\small\setlength{\parskip}{0.45em}}
  {\par\medskip\hrule\medskip}

% --- Mots-clés ---------------------------------------------------------------
%
% En anglais, en clôture du résumé de la lecture modernisée : le vocabulaire
% moderne sous lequel chercher ce que les feuillets construisent. Le manifeste
% du site (`npm run manifest`) les extrait pour étiqueter la cote entière —
% cette ligne est la source unique des tags, il n'y a pas de fichier à côté.
\newcommand{\keywords}[1]{\par\smallskip\noindent
  {\footnotesize\textsc{Keywords} --- \textit{#1}}\par}

% --- L'appareil critique ----------------------------------------------------

% Le numéro de feuillet, en marge. C'est l'ancre qui permet de régler un
% désaccord sur une formule en regardant *un* feuillet plutôt qu'une chemise
% de six cents. Le site s'en sert aussi pour tourner les pages du fac-similé
% quand on fait défiler la transcription.
\newcommand{\leaf}[1]{%
  \marginnote{\footnotesize\color{gray!70}\textsf{#1}}%
  \label{leaf:#1}\ignorespaces}

% Illisible : on ne devine pas. Un mot inventé qui se lit comme les autres est
% le pire résultat possible.
\newcommand{\ill}{\textcolor{red!60!black}{[\,\ldots\,]}}

% Lecture douteuse : le mot est donné, et signalé comme incertain.
\newcommand{\uncertain}[1]{\uline{#1}}

% Ajout de l'éditeur — une lettre manquante, un mot sous-entendu.
\newcommand{\add}[1]{\textcolor{blue!50!black}{[#1]}}

% Biffé par Grothendieck lui-même. Ce qu'il a rayé fait partie du document :
% une rature dit souvent où la pensée a changé de direction.
\newcommand{\struck}[1]{\sout{#1}}

% Note du transcripteur, sur l'état matériel du feuillet.
\newcommand{\note}[1]{\par\smallskip{\small\color{gray!60!black}\textit{[#1]}}\par\smallskip}

% Note marginale de Grothendieck, distincte du corps du texte.
\newcommand{\marginal}[1]{\par\smallskip\noindent\hspace{1em}%
  {\small\color{gray!60!black}#1}\par\smallskip}

% --- Titre ------------------------------------------------------------------

\AtBeginDocument{%
  \begin{center}
    {\large\bfseries Fonds Alexandre Grothendieck --- cote n\textsuperscript{o}~\grothFolder}\\[2pt]
    {\itshape\grothFoldertitle}\\[4pt]
    {\small feuillets \grothFirst--\grothLast\ (lot \grothBatch)}\\[2pt]
    {\footnotesize\color{gray!60!black}Transcription établie d'après le fac-similé numérisé
     par l'Université de Montpellier.\\
     Les lectures incertaines sont soulignées, les illisibles notées [\,\ldots\,],
     les ajouts entre crochets.}
  \end{center}
  \vspace{1em}}

\endinput


\folder{31}
\batch{4}
\pages{61}{80}
\dating{[1963-1973]}
\watermark{Édition de démonstration}
\foldertitle{SGA A [SGA 4] (Résidus de rédaction) : notes manuscrites (s.d.)}

\begin{document}

\section*{Fin de la proposition sur $H^{i}(U,F)$ et les recouvrements (page 61)}

\page{61}
\note{La page achève la Proposition commencée page 59 (batch 3), dont les
conditions sont numérotées (i), (i bis), (ii), (ii bis), (iii).}

Reste \struck{\ill{}} \uncertain{(ii)} $\Rightarrow$ (i bis). \ill{} ceci,
\ill{} \ill{} $F$ dans un \uncertain{injectif} \ill{} $C^{0}$ \ill{}
\[
\cdots \longrightarrow F \longrightarrow C^{0} \longrightarrow C^{1}
\longrightarrow \cdots
\]
L'hyp.\ (ii) implique que l'on a une suite exacte de \emph{flasques}, donc
$C^{0}$ étant flasque, il en est de même de $C^{1}$. Donc la résolution peut
être utilisée pour calculer la cohomologie, d'où $H^{i}(U,F) = 0$ pour
$i \geqslant 2$, etc.

\section*{cd des affines (pages 64 à 70)}
\note{Le titre est de sa main, porté sur la page 62, feuillet par ailleurs
vierge qui ouvre le développement suivant ; ce feuillet ne reçoit donc pas
de numéro de page.}

\page{64}
\textbf{Faisceaux $\ell$-élémentaires sur $X$} \quad [$X$ noeth.]

De la forme $i_{*}(G)_{U,X}$, i.e.
\begin{itemize}
\item $i : \operatorname{Spec} k(x) \longrightarrow X$ imm.\ canonique
  ($x \in X$)
\item $G$ un faisceau fini de $\ell$-torsion sur $X$
\item $U$ un \add{\uncertain{uniramifié}} ouvert de \struck{$X$} dans
  $Y = \bar{x}$, \quad $U = X - Z$
\end{itemize}
\note{Le mot inséré au-dessus de « ouvert » est d'une lecture douteuse ; le
$X$ qui suit « de » est barré d'une croix appuyée.}

\struck{Faisceau strictement élémentaire} \add{tel que} $G$ soit \emph{étale}
sur $U$, et $U$ est \emph{unibranche}.

\textbf{Proposition.} Tout faisceau de $\ell$-torsion noeth.\ sur $X$ a une
suite de composition finie des faisceaux de $\ell$-torsion élémentaires,
stricts.

\textbf{Cor 1.} Soit $n$ un entier. Conditions équivalentes :
\begin{enumerate}
\item[(i)] $\operatorname{cd}_{\ell}(X) \leqslant n$
\item[(ii)] Pour tout $Y$ intègre fermé dans $X$, \struck{et tout}
  \add{et pour $x$,} pour tout fermé $Z$ de $Y$, \struck{faisceau} tout
  $k(x)$-faisceau $G$ fini de $\ell$-torsion, si $i$ désigne l'immersion
  canonique $\operatorname{Spec} k(x) \to Y$, on a
  \begin{itemize}
  \item[a)] $H^{n}\bigl(Y, i_{*}(G)\bigr) \longrightarrow
    H^{n}\bigl(Z, i_{*}(G)|Z\bigr)$ surjectif
  \item[b)] $H^{n+1}\bigl(Y, i_{*}(G)\bigr) = 0$
  \end{itemize}
\end{enumerate}

\struck{De plus, on peut}

(ii bis) Rien ne change à (ii), sauf que $Y \to X$ est fini, \ill{} une
immersion, $Y$ intègre \emph{unibranche}, $i_{*}(G)$ \ill{} sur $X - Z$.

\textbf{Cor 1 bis.} Conditions équivalentes :
\begin{enumerate}
\item[(i)] Pour tout $F$ de $\ell$-torsion, et $T$ fermé,
  $H^{i}_{T}(X,F) = 0$ si $i > n$
\item[(ii)] \ill{}
  \begin{itemize}
  \item[a)] $H^{n}_{T}\bigl(Y, i_{*}(G)\bigr) \longrightarrow
    H^{n}_{T \cap Z}\bigl(Z, i_{*}(G)|Z\bigr)$ surjectif
  \item[b)] $H^{n+1}_{T}\bigl(Y, i_{*}(G)\bigr) = 0$
  \end{itemize}
\end{enumerate}
\note{Le bloc du Cor 1 bis est traversé d'un long trait courbe.}

\marginal{Si \ill{} une famille de $T$, \ill{} contenant $X$, \ill{}}

\page{65}
\textbf{Cor 2.} Soit $\varphi$ une fonction \struck{\uncertain{strictement}}
\emph{croissante} de \struck{\ill{}} $X$ dans $\mathbf{N}$.
\struck{Conditions} Pour toute partie fermée $Y$ de $X$, soit
\[
\varphi(Y) = \operatorname{Sup}_{x \in Y} \varphi(x)
 = \operatorname{Max}_{x \text{ maximal dans } Y} \varphi(x)
\]
Conditions équivalentes :
\begin{enumerate}
\item[(i)] Pour tout sous-préschéma fermé $Y$ de $X$,
  $\operatorname{cd}_{\ell}(Y) \leqslant \varphi(Y)$
\item[(ii)] Pour tout sous-préschéma fermé \emph{intègre} $Y$ de $X$, de
  \ill{} $n$, et tout faisceau $G$ fini de $\ell$-torsion sur
  $\operatorname{Spec} k(x)$, on a
  \[
  H^{n+1}\bigl(Y, i_{*}(G)\bigr) = 0 \qquad
  \text{si } n = \varphi(X) = \varphi(Y).
  \]
\end{enumerate}

(ii bis) Comme (ii), mais avec $Y$ fini sur $X$, intègre et unibranche.

Récurrence noethérienne \ill{} au cor.\ 1.

\struck{Cor 3} \textbf{Cor 2 bis.} Conditions équivalentes :
\begin{enumerate}
\item[(i)] Pour tout $F$ de $\ell$-torsion sur un sous-préschéma fermé $Y$ de
  $X$, et pour tout $T$ de $Y$, on a $H^{i}_{T}(Y,F) = 0$ si
  $i > \varphi(Y)$
\item[(ii)] \ill{} \quad $H^{n+1}_{T}\bigl(Y, i_{*}(G)\bigr) = 0$, \quad
  $n = \varphi(Y) = \varphi(x)$
\end{enumerate}
\note{Comme au bas de la page précédente, un long trait oblique traverse le
bloc.}

\page{66}
\struck{Lemme} \textbf{Prop.} Soit $\varphi$ une fonction sur $X$,
\add{\ill{}} satisfaisant les conditions suivantes :
\begin{itemize}
\item[a)] Pour tout $x \in X$, \quad
  $\varphi(x) \geqslant \operatorname{cd}_{\ell}\bigl(k(x)\bigr)$
\item[b)] Pour \add{$y \ne x$} $y \in \overline{\{x\}}$, \quad
  $\bar{y}$ \struck{\ill{}} une \ill{} géométrique sur \struck{$X$}
  \[
  \varphi(y) < \varphi(x) - \operatorname{cd}_{\ell}
  \bigl(\underline{O}_{X,\bar{y}} \otimes k(x)\bigr)
  \]
  \note{Sous la formule : « anneau des fonctions de
  $\underline{O}_{X,\bar y}$, \ill{} hensélisé strict ».}
\item[c)] Sous ces conditions, pour tout $F$ de $\ell$-torsion, on a
  \[
  H^{i}(X,F) = 0 \qquad \text{si } i > \operatorname{Sup}_{x \in
  \operatorname{Supp} F} \varphi(x) = \varphi(F).
  \]
\end{itemize}

\marginal{il y a intérêt à inclure le cas \uncertain{$\ell = 0$}, \ill{} du
faisceau \ill{} sur \ill{} ; dans b), n'y a-t-il \ill{} \ill{} que
\struck{\ill{}} devient \ill{} ; il y a intérêt à \uncertain{prendre}
$\varphi(x) = \dim$ de \struck{Zariski} \add{Krull} de $\bar{x}$}

\textbf{Dém.} \struck{\ill{}} Récurrence noeth.\ sur $X$. \struck{\ill{}}
\ill{} $F$ constructible, et comme \ill{}, il y a une \ill{},
$\varphi(F) = \operatorname{Sup}_{x \in Y} \varphi(x)$. \ill{}

Si $\operatorname{Supp} F \subset$ \struck{$\overline{\{x\}}$} \ill{}, alors
$H^{i}(Y,F) = 0$ si $i > \varphi(Y)$. \struck{\ill{}}

D'ailleurs, le cas \ill{} \ill{} $X$ intègre \ill{} fini. \ill{}, et $F$
\note{la phrase se poursuit page 67}

\page{67}
de la forme $g_{*}(G)$, i.e.\ $g : \struck{\operatorname{Spec}} \bar{x} \to
X$, $G$ fini sur \struck{$X$} $\bar{x}$, donc
$\varphi(F) = \varphi(\bar{X}) = \varphi(x)$
\[
H^{*}\bigl(\overline{\{x\}}, G\bigr) \Longleftarrow
H^{p}\bigl(X, R^{q} g_{*}(G)\bigr)
\]
\struck{Lemme. $R^{q} g_{*}(G)_{\bar{z}}$ \quad $E_{2}^{pq} = 0$ si
$p+q > \varphi(x)$ et $q > 0$.}

\add{Rq.} Pour $y \in X$, posons $\underline{O}_{X,\bar{y}} =$ anneau
strictement local \ill{} au-dessus de $y$ ou $\bar{y}$,
\[
K_{\bar{y}} = \underline{O}_{X,\bar{y}} \otimes_{X} k(x) = \ill{}
\]

\begin{tikzcd}[column sep=small, row sep=small]
x \arrow[d] & x_{\bar y} \arrow[l] \arrow[d] \\
X & X_{\bar y} \arrow[l]
\end{tikzcd}

\textbf{Rq.} \struck{\ill{}} $R^{q} g_{*}(G)_{\bar y} =
H^{q}\bigl(K_{\bar y}, G|K_{\bar y}\bigr)$, donc \ill{}
$q \leqslant \operatorname{cd}_{\ell}(K_{\bar y})$, donc \ill{} satisfait
\[
\varphi(y) < \varphi(x) - \operatorname{cd}_{\ell}(K_{\bar y})
 \leqslant \varphi(x) - q
\]
d'où
\[
\boxed{\ \varphi\bigl(R^{q} g_{*}(G)\bigr) \leqslant \varphi(x) - q\ }
\]
D'autre part, si $q > 0$, \ill{} $R^{q} g_{*}(F)_{x} = 0$, et l'hyp.\ de
récurrence \ill{} :
\[
H^{p}\bigl(X, R^{q} g_{*}(G)\bigr) = 0
\]
\note{Un croquis du plan $(p,q)$ occupe la marge droite : des croix
au-dessous de la diagonale, et la droite $p+q = \varphi(x)$ marquée
« $\varphi(x)$ » et « $\varphi(x)+1$ » sur l'axe des $p$ ; la région
$p+q \geqslant \varphi(x)$ ne rencontre aucune croix.}

\page{68}
\textbf{N.B.} Si $f : X \to Y$ morphisme de type fini de \ill{} \ill{}, pour
$x \in X$, $y = f(x)$,
\[
\varphi(x) \leqslant \varphi(y) + 2 \deg\!\operatorname{tr}
\bigl(k(x)/k(y)\bigr)
\]

\page{69}
\textbf{Corollaire.} Supposons que pour tout \add{noeth.\ \ill{}} \ill{} $Y$
de $X$, et tout $y \in X$ \ill{} $\bar{y}$ de $Y$, les anneaux \ill{}
$K_{x,\bar y}$, $\underline{O}_{X,\bar y}$ \struck{soient de}
$\operatorname{cd}_{\ell}$ satisfont
\[
\operatorname{cd}_{\ell} K_{x,\bar y} \leqslant \dim \underline{O}_{X,\bar y}
\quad \bigl[= \dim \underline{O}_{Y,y}\bigr].
\]
\struck{Alors} \ill{} \struck{\ill{}}

Soit
\[
\boxed{\ \varphi(x) = \operatorname{Sup}_{y \in \bar{x}}
\Bigl(\operatorname{cd}_{\ell}\bigl(k(y)\bigr) + 2 \dim
\underline{O}_{\bar{x},y}\Bigr)\ }
\]
Alors, pour toute partie \ill{},
\[
\varphi(Y) = \operatorname{Sup}_{y \in Y}\Bigl(
\operatorname{cd}_{\ell}\bigl(k(y)\bigr) + 2 \dim
\underline{O}_{Y,y}\Bigr), \qquad
\operatorname{cd}_{\ell}(X) \leqslant \varphi(X).
\]

\marginal{\textbf{Question.} $A$ local noeth.\ hensélien, \add{intègre, de
\ill{}} \ill{} résidu $k$ de \struck{dimension} \ill{}, tel que
$\operatorname{cd}_{\ell}(k) < +\infty$. Alors,
$\operatorname{cd}_{\ell} K \geqslant \dim A + \operatorname{cd}_{\ell} k$ ??}

\textbf{Exemples.}
\begin{enumerate}
\item[1)] Schémas algébriques, plus généralement \ill{} \ill{} de type fini.
  [N.B. \ill{} \ill{} du degré de transcendance et du degré de transcendance
  dans \ill{} de dimension \ill{}]
\item[2)] Corps de \ill{} \ill{} nombres
\item[3)] \uncertain{Compactification} \ill{} \quad $A$ strictement local
  \ill{}, donné le $\operatorname{Spec} A$
\item[4)] Schémas \ill{} \emph{affine de dim 1}
\end{enumerate}

\page{70}
\struck{$R^{q} i_{*}(G|F)$}
\[
H^{*}(G) \Longleftarrow E_{2}^{pq} = H^{p}\bigl(X, R^{q} i_{*}(G)\bigr)
\]
\note{Un croquis du plan $(p,q)$ suit, la diagonale marquée « $n$ » et
« $n+1$ » sur les deux axes.}

$x \in \operatorname{Supp} R^{q} i_{*}(G)$ \qquad
$\operatorname{cd}(\bar{x}) + \varphi(x)$ \qquad
$\varphi$ \quad $\operatorname{cd}(\bar{x}) + q \leqslant$
\[
\underline{O}_{X,\bar y} \otimes_{X} k(x), \qquad
q \leqslant \operatorname{cd}\bigl(\underline{O}_{X,\bar y} \otimes \ill{}
\bigr), \qquad
q \leqslant \operatorname{cd}_{p}\operatorname{Frac}
\underline{O}_{X,\bar y}
\]
\[
\operatorname{cd}(\bar{y}) + q \ \operatorname{cd}_{p}\operatorname{Frac}
\underline{O}_{X,\bar y} < n
\]
\struck{$\varphi(y) < \varphi(x) - \operatorname{cd}_{\ell} \ill{}$}
\[
H^{*}(G) \longrightarrow H^{*}f' \longrightarrow H^{k+1}(X_{T})
\longrightarrow H^{k+1}(G)
\]

\marginal{$d'$ \quad ; \quad $2d' + \operatorname{cd}(\ ) < 2d+1$ \quad ;
\quad $(d - d' + 1)$ \quad ; \quad $d + d' + 1 \leqslant 2d+1$ \quad ; \quad
$2d+1$ \quad ; \quad \struck{$2d$} $\leqslant d+1$ \quad ; \quad
$1 + 2d' + (d - d')$ \struck{$2(d-1)$} $= d + d' + 1 \leqslant 2d < 2d+1$}

$\operatorname{cd}_{\ell}(X) \leqslant n$

\textbf{Lemme.}
\[
\boxed{
\begin{array}{l}
\operatorname{cd}_{\ell}\bigl(k(x)\bigr) \leqslant n \\[2pt]
\forall\, Y \ \uncertain{\text{réunion}} \text{ dans } X, \quad
\operatorname{cd}_{\ell}(Y) < n \\[2pt]
\forall\, y \in X,\ y \ne x, \quad
\operatorname{cd}_{\ell}(\bar{y}) + \operatorname{cd}_{\ell}
\operatorname{Frac} \underline{O}_{X,\bar y} < n
\end{array}}
\]
\[
2 \dim \bar{y} + \dim \underline{O}_{X,\bar y} \quad
\bigl[\oplus \dim X - \dim \bar{y}\bigr] \qquad
\dim X + \dim \bar{y} \leqslant 2 \dim X
\]

\section*{Feuillet de tapuscrit corrigé : quotients par un groupe de type multiplicatif, 2.3 à 2.5 (page 71)}
\note{Le feuillet est dactylographié ; les références laissées en blanc sont
complétées à l'encre de sa main, les « VII » corrigés en « VIII », et une
parenthèse entière est biffée. Insertions dans \texttt{add}, suppressions
dans \texttt{struck}.}

\page{71}
morphismes fidèlement plats et quasi-compacts (cf.\ Exp IV, \add{3.4.}), de
plus $Y = X/G$ est affine sur $S$. Si de plus $X$ est de présentation finie
(resp.\ de type fini) sur $S$, il en est de même de $Y$.

La première assertion résulte de Exp \add{VIII} 5.1., \struck{et de} traitant
le cas où $G$ est diagonalisable, et de Exp IV, \add{3.5.2.}, qui permet de
s'y ramener, compte tenu que les morphismes fidèlement plats et
quasi-compacts \add{$S' \to S$} sont des morphismes de descente
\emph{effective} pour la catégorie fibrée des schémas affines sur d'autres,
i.e.\ pour tout $Y'$ affine sur $S'$, muni d'une donnée de descente
relativement à $S' \to S$, cette donnée de descente est effective, i.e.\ $Y'$
provient d'un $Y$ affine sur $S$, cf.\ SGA VIII 2.1. Pour la deuxième
assertion, on est ramené également au cas diagonalisable\add{, Exp VIII
5.8.}, car les conditions de finitude envisagées se descendent par morphismes
fidèlement plats et quasi-compacts \struck{(pour le type fini, c'est dans
dans ce cas, reprenant les notations de Exp VI, 5.1., on aura
$X = \operatorname{Spec}(A)$, $X = \operatorname{Spec}(A_{0})$, où $A_{0}$
est l'Algèbre formée \add{la} des composantes de degré 0 de l'Algèbre graduée
$A$)} (SGA VIII 3.3. et 3.6.).

Procédant comme dans Exp \add{VIII}, corollaires 5.5. à 5.7., on tire de 2.3.\ :

\textbf{Corollaire 2.4.} Sous les conditions de 2.3., le morphisme graphe
\[
X \times_{S} G \longrightarrow X \times_{S} X
\]
est une immersion fermée. Pour toute section $s$ de $X$ sur $S$, le morphisme
correspondant $g \mapsto s \cdot g$ :
\[
G \longrightarrow X
\]
est une immersion fermée.

En particulier :

\textbf{Corollaire 2.5.} Soit $u : G \to H$ un monomorphisme de
$S$-préschémas en groupes, avec $G$ de type multiplicatif et $H$ affine sur
$S$. Alors $u$ est une immersion fermée, \struck{et} $H/G = Y$ existe et est
affine sur $S$, enfin $H$ est un fibré principal homogène sur $Y$ de groupe
$G_{Y}$.

\section*{Deux théorèmes de finitude pour $\operatorname{cd}_{\ell}$ sur un anneau strictement local (pages 73 à 80)}

\page{73}
\textbf{Théorème 1.} Soit $A$ anneau str.\ local noeth.\ de dim $n$,
\marginal{(strictement) str.\ quasi-\ill{} de \ill{} stricte}
$Y = \operatorname{Spec}(B)$, \struck{alors pour tout} $F$ un faisceau de
$\ell$-torsion sur $X = \operatorname{Spec} A$ \add{$\ell$ premier à la car.\
résiduelle} tel que
\[
x \in \operatorname{Supp} F \ \Longrightarrow\ \dim \underline{O}_{X,x}
\geqslant m
\]
(i.e.\ \ill{} ; $F$ est \add{« nul en codim $< m$ »}).
\struck{Dans le cas où le \ill{} montre \ill{}}

Alors $H^{i}(X,F) = 0$ pour $i > n + 1 - m$ \struck{\ill{} de $X_{0}$.}

[En particulier, $\operatorname{cd}_{\ell}(X) \leqslant n+1$, et pour tout
fermé $Z$, $\operatorname{cd}_{\ell}(Z) \leqslant n + 1 -
\operatorname{codim}(Z,X)$.]

\textbf{Théorème 2$_{n}$.} Soit $A$ \ill{} th.\ local noeth.\ de dim $n$, et
\struck{\ill{}} $f \in A$, $X = \operatorname{Spec}(A)$, \struck{$M$} $=
X_{f}$, alors
\[
\operatorname{cd}_{\ell}(X_{f}) \leqslant n
\]
pour tout $\ell$ premier à la car.\ résiduelle.

[N.B. $1_{n} \Rightarrow 2_{n}$ car \ill{} (cas \ill{} $= 1_{n+1}$ \ill{})]

\textbf{Lemme 1.} $2_{n'}$ pour $n' \leqslant n+1$ $\Rightarrow$ $1_{n}$
\quad \ill{}
\note{Sous « $n' \leqslant n+1$ », une accolade portant « $2'_{n+1}$ ».}

\ill{} pour \ill{} $Y \subset X$ \ill{} de codim $\geqslant 1$,
b) ne \ill{}\note{« b) ne \ill{} » est cerclé.} \ill{} Alors $\operatorname{cd}_{\ell}(S) \leqslant
n+1-m$.

En effet, soit $Y \subset X$ fermé \ill{} de codim $m$, soit
$\bar{X} = \mathbf{P}^{1}_{Y}$, considérons $\bar{S}$ l'adhérence de $S$,
considérons $T = \bar{S} - S$, c'est un schéma str.\ local
($\subset \bar{X} - x$), où $x$ est le pt fermé. Soit $F$ de $\ell$-torsion
sur $S$ :
\[
\cdots \longrightarrow H^{i}(\bar{S},F) \longrightarrow H^{i}(S,F)
\longrightarrow H^{i+1}_{T}(\bar{S},F)
\]
Or
\[
H^{i}(\bar{S},F) = H^{i}(\bar{S}_{0},F) = H^{i}(\mathbf{P}^{1}_{k}, F_{0}),
\quad (\text{nul si } i \geqslant 3)
\]
\[
H^{i+1}_{T}(\bar{S},F) = H^{i}\bigl(\bar{S}_{x} - T_{x}, F_{(x)}\bigr)
\quad \text{si } i \geqslant 2
\]
D'ailleurs $\dim \bar{S}_{x} \leqslant n+1-m$ \add{\ill{}} \ill{}

\page{74}
\struck{car en localisant en \ill{} géométriques}
\[
\dim \bar{S}_{x} + \operatorname{codim}\bigl(\bar{S}_{x}, \bar{X}_{x}\bigr)
\leqslant \dim \bar{X}_{x} = n+1
\]
\note{Une accolade sous le terme de codimension porte « $\geqslant m$ ».}

Dans \struck{\ill{}} l'hyp.\ de récurrence $2_{n'}$ pour $n' \leqslant n+1$
[en particulier $2_{n+1-m}$], on en conclut
\[
H^{i+1}_{T}(\bar{S},F) = 0 \qquad \text{si } i > n+1-m
\]
Si $n+1-m \geqslant 2$, i.e.\ $m \leqslant n-1$, \ill{} on a
\[
H^{i}(\bar{S},F) = H^{i}(\bar{S}_{0},F_{0}) = 0 \quad \text{si }
i > n+1-m \geqslant 2
\]
d'où $H^{i}(S,F) = 0$. Reste le cas \ill{} \ill{} ($S$ \ill{} $\ne
\emptyset$) \ill{} \struck{si $m > n+1$}, \ill{} $m \geqslant n$,
\struck{si $m = n+1$}, \struck{\ill{}} \ill{} $X_{0}$, $A$ \ill{}
$H^{i}(S,F) = 0$ \ill{} pour $i > 0 = (n+1)-m$, ok.

Si $m = n$, alors \ill{} de dim $\leqslant 1$. \struck{Par la}
\add{si $S \ne x_{0}$} \struck{\ill{}} \ill{} $= X_{0}$, \ill{}. $\bar{S}$
\ill{} de dim 0, donc $H^{i}(\bar{S}_{0},F) = 0$ si $i > 0$ \ill{} et
$i > n+1-m = 1$. Pour $S = X_{0}$, \ill{} \add{\ill{}} \ill{} sup.\ de
\ill{}]. \ill{} $k(x)$ \ill{} $\dim$ \ill{} \emph{cd} $<
\operatorname{cd}_{\ell} \leqslant 1$, \ill{} qui \ill{}

\page{75}
\textbf{Lemme 2.} Se borner \add{aux anneaux \ill{} sur un corps}, i.e.\
\ill{}
\[
\struck{1'_{n} \Rightarrow 2'_{n+2}} \qquad
1'_{n} + 2'_{n+1} \Longrightarrow 2'_{n+2}
\]
[N.B. Et $2'_{n+2} \Rightarrow 1'_{n+1} + 2'_{n+2}$ qui \ill{} 1, OK]

En effet, \ill{} \ill{} (Zariski \ill{}). Et \struck{la dernière \ill{}}
$f \in \mathfrak{m}_{A}$ \add{$2'_{n+2}$} \ill{} (\ill{} $X_{f} =
\emptyset$ et \ill{} trivial]. \struck{Alors $f$ de \ill{}} \ill{} les \ill{}
$A$ intègre \ill{} $f$ \ill{} \ill{} $i > 2$ \ill{}
\note{Deux croquis --- un cube et un parallélogramme --- occupent la marge
gauche.}

\struck{$A_{1} = k\langle x_{1}, \ldots, x_{m}\rangle \to A$} \quad
$\to k\langle x_{1}, \ldots, x_{n+2}\rangle \to A$, et $X =
\struck{\operatorname{Spec}}$ \ill{} \quad $X' = \operatorname{Spec} A'$,
\struck{$M$} \ill{} \struck{$M'$} $= X'_{f}$. \ill{}
$\operatorname{cd}_{\ell}(X') \leqslant n$, donc \ill{}
\[
k\langle x_{1}, \ldots, x_{n+2}\rangle, \quad f = x_{1}\ldots \ill{}
\]
qui est de dim $n+1$, et $B = k\langle x_{1}, \ldots, x_{n+1}\rangle$ \ill{}
$Y = \operatorname{Spec} B$, $U = Y_{f}$ \ill{} $V \to U$. Pour l'hyp.\
\ill{} prouver que $f$ \ill{} $R^{q} f_{*}(F)$ \ill{} codim
$i \leqslant q-1$, car alors, considérons le \uncertain{prolongement}
$2'_{n+1}$\note{$2'_{n+1}$ est cerclé.}, il suffit de \ill{} $R^{q} f_{*}(F)$ \ill{} nul en
\ill{}
\[
H^{p}\bigl(U, R^{q} f_{*}(F)\bigr) \ne 0 \ \Longrightarrow\ \struck{\ill{}}
\ p \leqslant n+1-(q-1), \ \text{i.e.\ } p+q \leqslant n+2,
\]
d'où $H^{i}(V,F) \ne 0 \Rightarrow i \leqslant n+2$.

\marginal{Attention : conditions \ill{} \ill{} ?}

\page{76}
On les \ill{} locaux \ill{} de $U$ \ill{} de dim $i \leqslant n$ ! Donc
\ill{} $=$ les \ill{} stricts locaux $B'$ de $U$, et \struck{de} $V$ \ill{}
\ill{} $H^{i}$ \ill{} les dim.\ \emph{relatives}
\[
V' = \struck{V} \otimes_{B} B' \quad \add{\text{Affine}} \ill{}
\]
\ill{} relative 1 sur $U$, donc \ill{} dim relative 1 sur $B'$, \ill{}
consort par l'hyp.\ $1'_{n}$\note{$1'_{n}$ est cerclé.} \ill{} indiquée car
\[
H^{q}(V',F') = 0 \quad \text{si } q > i+1, \ \text{i.e.\ }
\underline{i < q-1},
\]
[i.e.\ $R^{q} f_{*}(F)$ \ill{} nul en codim $i < q-1$].

C'est ce qu'il \ill{} \struck{fallait vérifier}, rien \ill{} (sauf pour les
cas t.f.\ \ill{} conclus]

\textbf{Lemme 3.} $1'_{n}$ \struck{\ill{}} (implique que \ill{} \ill{} affine
$X$, \ill{} de t.f.\ sur \ill{} corps sép.\ clos $k$, \ill{} de dim de Krull
$d \leqslant n+1$, \ill{}
\[
\operatorname{cd}_{\ell}(X) \leqslant \dim X
\]

\textbf{Dém.} \ill{} \ill{} $X$ \ill{}, Trivial si $d \leqslant 0$. \ill{}
$\operatorname{cd}(n-1)$ \ill{}

\struck{Récurrence sur $d$} \quad Récurrence sur $d = \dim X$, $d \leqslant
n+1$

\struck{Par le \ill{}} tout $d \geqslant 1$. \ill{} $X = E^{d}_{k}$.
\[
\operatorname{cd}_{d}\bigl(E^{d}_{k}\bigr) \leqslant d
\]
\struck{$d \leqslant n+1$} \quad Trivial si $d \leqslant 0$.

En \ill{} $X = E^{d}_{k} \xrightarrow{\ f\ } Y = E^{d-1}$, \ill{}
$d' = d-1 \leqslant n$. En vertu de $1'_{n}$,
\[
R^{q} f_{*}(F)_{\bar y} \ne 0 \ \Longrightarrow\
\dim \underline{O}_{Y,y} + 1 = \struck{i+1} \ \text{i.e.\ }
R^{q} f_{*}(F)_{\bar y} \ \ill{}
\]

\page{77}
« nul en codim $\leqslant q-1$ ». Par l'hyp.\ de récurrence il s'ensuit que
\[
H^{p}\bigl(Y, R^{q} f_{*}(F)\bigr) \ne 0 \ \Longrightarrow\
p \leqslant n - (q-1), \quad \text{i.e.\ } p+q \leqslant n+1.
\]
\marginal{$n-q+1$}
Donc $E_{2}^{pq} = 0$ \ill{} $p+q \leqslant n+1$, d'où la conclusion \ill{}

Plus \ill{}, les Th.\ 1 et 2 sont prouvés dans la restriction
« \struck{$p = $ t.f.\ est} \add{$\mathbf{E}$-ponctuellement} t.f.\ » \ill{}
pour « \ill{} de t.f.\ ». N.B. Mais, \ill{}

Th.\ 1 et 2 sont prouvés « \struck{loc.\ t.f.\ sur} sur $k$ ». Il en est
ainsi de Lemme 3, \ill{} t.f.\ sur $k$ ».

\ill{} maintenant le \ill{} il s'agit de prolongement, \ill{} de t.f.\ \ill{}
de dim $\leqslant 1$. \ill{} le \struck{\ill{} local $2_{n}$}, \ill{}, bien
si $n = 0$, \ill{} $n > 0$, et le th.\ \ill{} \ill{} $A$ str.\ local
$1 \ne A$ local \struck{\ill{}} $A$,

\marginal{\ill{}}

\page{78}
\note{Un croquis ouvre la page : une sphère dont l'équateur porte $V(f)$ et
qu'un petit ovale vertical, $V(\pi)$, rencontre ; à gauche, $X \to Y$ et une
accolade marquée $V$.}

Soit $\pi \in \mathfrak{m}_{A}$, $\pi \ne 0$, \ill{} $\pi$ \ill{} div.\ de 0.
Alors \ill{}, \struck{$H^{i}$ \ill{} $V = X_{f}$} \add{$F$}
\[
V = X_{f}, \qquad U = X_{f\pi} = V_{\pi}, \qquad
Z = V - U = V \cap V(\pi) = V(\pi)_{f}
\]
\[
H^{i}_{Z}(V,F) \longrightarrow H^{i}(V,F) \longrightarrow H^{i}(U,F)
\]
et il \ill{} \ill{} pour les extrémités \ill{} \ill{}. On \ill{} corps des
fractions, $K$ de $A$, \ill{} $K$ est de dim tot.\ $\leqslant 1$, et \ill{}
de $\ell$-dim ch.\ $\leqslant (n-1)+1 = n$, \ill{} qui \ill{} des dernières.
\[
H^{*}_{Z}(V,F) \Longleftarrow H^{p}\bigl(Z, \underline{H}^{q}_{Z}(F)\bigr)
\]
comme $Z = V(\pi)_{f}$ \add{\ill{} dim $h(A)$}, \ill{} dim comb.\ de $Z$ qui
est \ill{} son \ill{}-schéma \ill{}. $E_{2}^{pq} = 0$ pour $q \ne 0,1$ \ill{}
donne \ill{} $H^{i}_{Z}(V,F)$ pour $i \geqslant n+1$.

\page{79}
\[
\underline{H}^{q}_{Z}(F)_{\bar x} \ne 0 \ \Longrightarrow\
\dim \underline{O}_{Z,x} \geqslant q-1
\]
i.e.\ $\underline{H}^{q}_{Z}(F)_{\bar x}$ est nul en codim $< q-1$,

[ce qui indique que
\[
H^{p}\bigl(Z, \underline{H}^{q}_{Z}(F)\bigr) \ne 0 \ \Longrightarrow\
p \leqslant (n-1)-(q-1) = n-q, \quad \text{i.e.\ } p+q \leqslant n,
\]
d'où la conclusion voulue]. Il faut donc \ill{} que
\[
q-1 > \dim \underline{O}_{Z,x} \ \Longrightarrow\
\underline{H}^{q}_{Z}(F)_{\bar x} = 0.
\]
\ill{} évidemment \ill{}. Or, pour $q = 0,1$,
\[
\underline{H}^{q}_{Z}(F)_{\bar x} = H^{q-1}\bigl(X_{\bar x} -
V(\pi)_{\bar x}, F\bigr)
\]
qui est \ill{} \ill{} si $q-1 > \dim X_{\bar x}$ \ill{} [car
$X_{\bar x}$ \ill{} \ill{} t.f.\ \ill{} un corps sép.\ clos $K$]

\textbf{Théorème} \struck{Question}. Soit $A$ str.\ local, $X$ de type fini
et affine, $Y = \operatorname{Spec} A$, $Y' = Y - \{\text{pt}\}$,
$X' = X \times_{Y} Y'$, \struck{la dimension de Krull}, soit
\struck{$X_{0}$}
\[
d = \operatorname{Sup}\bigl[\dim X_{0},\ \dim X' + 1\bigr].
\]
Alors on \ill{} pour $H$ \ill{} \ill{} de $A$, \ill{}
$\operatorname{cd}_{\ell}(X) \leqslant d$
\[
\operatorname{cd}_{\ell}(X) \leqslant \operatorname{Sup}
\bigl[2 \dim X_{0},\ 2 \dim X' + 1\bigr]
\]
\[
\operatorname{cd}_{\ell}(X) \leqslant 2 \dim Y - 1 + 2d \qquad
\bigl(\text{si } \underline{\dim Y \ne 0}\bigr)
\]
\marginal{N.B. Si $X$ est intègre et domine $Y = \operatorname{Spec} A$ [ce
qui n'est pas trop restrictif] on trouve
$\operatorname{cd}(X) \leqslant \dim Y + d$, « $d = \dim$ de la fibre
générique » (cerclé), resp.\ ($X$ irréd.\ dominant)}

\page{80}
\begin{tikzcd}[row sep=small]
E = Y[t_{1}, \ldots, t_{r}] \arrow[d] \\
Y
\end{tikzcd}

\ill{} \add{de dim $\varphi$)} \ill{} plongeons $Y$ intègre, $X$ \ill{} dans
$E = Y[t_{1}, \ldots, t_{r}]$, $X$ intègre. Alors \struck{codim} \ill{} pour
\[
\operatorname{codim}(X,E) \geqslant n+r-\nu \quad \text{i.e.\ }
\nu \geqslant \boxed{n+r-\operatorname{codim}(X,E)}
\]
$E$ \ill{} affine, si $X \ne X_{0}$, alors le premier \ill{} égal à
\[
\operatorname{codim}(X',E') = \dim E' - \dim X' = \bigl((n-1)+r\bigr) -
\dim \ill{}
\]
\[
\dim X' \geqslant (n-1+r) - (\nu-1) = n+r-\nu \qquad \text{ok}
\]
Si $X = X_{0}$, alors \struck{\ill{} dim ou $X_{0}$}
\[
\operatorname{codim}(X,E) = \operatorname{codim}(X_{0},E_{0}) +
\operatorname{codim}(E_{0},E) = (r-\nu) + n \qquad \text{ok.}
\]
[N.B. On \ill{} \ill{} $\operatorname{codim}(X,E) > n+r-\nu$,
\struck{est-ce vraiment ?} \ill{} lorsque $\dim X_{0} > \dim X' + 1$ \ill{}
$X \ne X_{0}$.]
\[
\operatorname{cd}_{\ell}(X) \leqslant (n+r) - \operatorname{codim}(X,E)
\]
\note{Le membre de droite est raturé et réécrit à plusieurs reprises.}

On ne \ill{} \ill{} Pour $r = 0$, c'est trivial (alors $X \subset Y$, \ill{}
$H^{i}(X,F) = 0$ pour $i > 0$ i.e.\ \ill{}
$n \geqslant \operatorname{codim}(X,Y) = \dim X$. Si
$\operatorname{cd}(X) = 0 \leqslant$ \struck{\ill{}} \ill{} th.\ $1_{m}$.
Dans le cas $r = 1$, c'est le th.\ \ill{} \struck{gratuit} \add{$r \geqslant
2$}, \ill{} \ill{} spéciale \ill{}
\[
E^{r}_{Y} \longrightarrow E^{r-1}_{Y} \qquad \cdots
\]

\end{document}
